% Part: proof-theory
% Chapter: sequent-calculus
% Section: introduction

\documentclass[../../../include/open-logic-section]{subfiles}

\begin{document}

\olfileid{pt}{seq}{int}

\olsection{ভূমিকা}

সিকোয়েন্ট কলন হলো !!{proof} ব্যবস্থা, যা
১৯৩০-এর দশকে গেরহার্ড জেন্টজেন প্রবর্তন করেন।
কিছু !!{proving} করার ব্যবহারিক পদ্ধতি দেওয়া
এর মূল উদ্দেশ্য ছিল না। বরং জেন্টজেন
এগুলি ব্যবহার করেছিলেন !!{proof}s-এর সম্ভাব্য
গঠন ও বৈশিষ্ট্য সম্বন্ধে কয়েকটি ফল প্রমাণের
জন্য। প্রমাণতত্ত্ব ঠিক এই ধরনের প্রশ্ন নিয়েই
আলোচনা করে।

নাম থেকেই বোঝা যায়, সিকোয়েন্ট কলনে
!!{formula}s নয়, \emph{সিকোয়েন্ট} নিয়ে
কাজ হয়। একটি সিকোয়েন্ট হলো !!{formula}s-এর
দুটি অনুক্রম বা বহুসেটের যুগল। জেন্টজেনের
মূল ব্যবস্থাগুলি (\Log{LK} ও \Log{LJ})
অনুক্রম ব্যবহার করত; এখন প্রমাণতাত্ত্বিকেরা
বহুসেট ব্যবহার করতে বেশি পছন্দ করেন।
বহুসেট সেটের মতো, তবে তাতে একই
!!{element} একাধিকবার থাকতে পারে।
অবশ্য সেটের মতো এখানেও !!{element}s-এর
ক্রম গুরুত্বপূর্ণ নয়। তাই
$\{!A, !A, !B\}$ ও $\{!A, !B, !A\}$
একই বহুসেট বোঝায়। কিন্তু সেই
বহুসেট $\{!A, !B\}$ এবং
$\{!A, !A, !B, !B\}$ থেকে আলাদা;
শেষের দুটিতে $!A$ ও $!B$-এর
পুনরাবৃত্তি-সংখ্যা প্রথমটির থেকে আলাদা।

প্রমাণতত্ত্বে রীতি অনুসারে !!{formula}s-এর
বহুসেট বড় হাতের গ্রিক অক্ষরে লেখা হয়।
তাই $\Gamma$, $\Delta$, $\Pi$,
$\Lambda$ বা $\Theta$ লিখলে আমরা
যে যুক্তিবিদ্যা নিয়ে কাজ করছি,
তার !!{formula}s-এর বহুসেট বুঝি।
আবার রীতি অনুসারে
$\tuple{\Gamma, \Delta}$ না লিখে
প্রমাণতাত্ত্বিকেরা দুই অংশের মাঝখানে
সিকোয়েন্ট-চিহ্ন বসান; আমরা
$\Sequent$ ব্যবহার করি।

\begin{defn}[সিকোয়েন্ট]
\emph{সিকোয়েন্ট} হলো এই রূপের রাশি:
\[
\Gamma \Sequent \Delta
\]
এখানে $\Gamma$ ও $\Delta$ হলো
!!{formula}s-এর সসীম (সম্ভবত ফাঁকা)
বহুসেট। $\Gamma$-কে \emph{পূর্বাংশ}
এবং $\Delta$-কে \emph{উত্তরাংশ} বলে।
\end{defn}

ধরি, $!G$ হলো~$\Gamma$-র
!!{formula}s-এর যৌগ
(ফাঁকা যৌগকে~$\ltrue$ ধরা হচ্ছে),
আর $!D$ হলো~$\Delta$-র
!!{formula}s-এর বিযোজন
(ফাঁকা বিযোজনকে~$\lfalse$ ধরা হচ্ছে)।
তাহলে কোনো কাঠামো~$\Struct{M}$-তে
সিকোয়েন্ট $\Gamma \Sequent \Delta$
তখনই সত্য, যখন $!G \lif !D$ সত্য।
চিরায়ত যুক্তিতে এর অর্থ,
হয়~$\Gamma$-র !!{formula}s-এর
অন্তত একটি মিথ্যা,
নয়তো~$\Delta$-র !!{formula}s-এর
অন্তত একটি সত্য।

$\Gamma$ যদি !!{formula}s-এর বহুসেট হয়,
তবে $\Gamma$-তে $!A$ যোগ করে পাওয়া
বহুসেটকে $\Gamma, !A$ (অথবা
$!A, \Gamma$) লিখি। $\Delta$-ও
যদি !!{formula}s-এর বহুসেট হয়,
তাহলে $\Gamma, \Delta$ হলো দুই
বহুসেটের বহুসেটীয় সংযুক্তি:
% BN-SRC-904: context-sensitive Bengali multiset wording.
$\Gamma$-তে $!A$-এর
\emph{পুনরাবৃত্তি-সংখ্যা}~n হলে
(অর্থাৎ $!A$-এর $n$টি অনুলিপি থাকলে)
এবং $\Delta$-তে সেই সংখ্যা~$m$
হলে $\Gamma, \Delta$-তে $!A$-এর
পুনরাবৃত্তি-সংখ্যা~$n+m$।
সিকোয়েন্টের কোনো এক পাশের
বহুসেট ফাঁকা হলে সাধারণত
সেই জায়গা ফাঁকাই রাখি।
যেমন, $\Sequent !A$ সিকোয়েন্টের
পূর্বাংশ ফাঁকা, আর উত্তরাংশে
$!A$-এর পুনরাবৃত্তি-সংখ্যা~$1$।

সিকোয়েন্ট কলনে !!^a{proof} হলো
সিকোয়েন্টের একটি বৃক্ষ। তার
একেবারে উপরের (বা প্রারম্ভিক)
সিকোয়েন্টগুলি আলোচ্য ব্যবস্থার
স্বতঃসিদ্ধ। কোনো সিকোয়েন্ট
এক বা দুটি অন্য সিকোয়েন্টের
নিচে থাকলে সেটি অবশ্যই
ব্যবস্থাটির কোনো অনুমান-বিধি
অনুসারে সেগুলি থেকে সঠিকভাবে
আসতে হবে। প্রথমে চিরায়ত
যুক্তিবিদ্যার দুটি সিকোয়েন্ট
ব্যবস্থা, $\Log{G1c}$ ও~$\Log{G3c}$,
দেখব। তাদের স্বতঃসিদ্ধ ও বিধি
\cref{pt:seq:int:tab:G1c,pt:seq:int:tab:G3c}-এ
দেওয়া আছে। \oliflabeldef{fol:seq::chap}{ (জেন্টজেনের
মূল ব্যবস্থা \Log{LK}-র বিবরণ
\olref[fol][seq][]{chap}-এ আছে।)}{}

\subfile{rules-G1c}
\subfile{rules-G3c}

ব্যবস্থা দুটির মধ্যে সূক্ষ্ম পার্থক্য
আছে; ফলে তাদের প্রমাণতাত্ত্বিক
বৈশিষ্ট্যও আলাদা। \Log{G1c}-র
স্বতঃসিদ্ধ $!A \Sequent !A$;
এতে অন্য কোনো !!{formula}s
থাকতে পারে না। \Log{G3c}-র
স্বতঃসিদ্ধ $!C, \Gamma \Sequent
\Delta, !C$ রূপের। এতে
$\Gamma$ ও~$\Delta$-তে
যেকোনো “পার্শ্ববর্তী”
!!{formula}s থাকতে পারে,
কিন্তু দুই পাশে থাকা
!!{formula}~$!C$ অবশ্যই
\emph{আণবিক} হতে হবে।
\Log{G1c}-তে $\LeftR{\land}$
ও $\RightR{\lor}$ বিধির
দুটি করে রূপ আছে; \Log{G3c}-তে
প্রতিটির একটি করে। এছাড়া
\Log{G1c}-তে “গঠনগত বিধি”
আছে: “সংকোচন”
(\LeftR{\Contraction}, \RightR{\Contraction})
ও “দুর্বলীকরণ”
(\LeftR{\Weakening}, \RightR{\Weakening})।

\Log{G1c} ও \Log{G3c}
উভয় ব্যবস্থাই চিরায়ত
যুক্তির জন্য বিশুদ্ধ ও পূর্ণ।
অর্থাৎ এই ব্যবস্থাগুলিতে
!!{provable} প্রতিটি সিকোয়েন্ট
প্রত্যেক !!{structure}-এ সত্য;
আবার প্রত্যেক !!{structure}-এ
সত্য এমন প্রতিটি সিকোয়েন্টের
!!a{proof} আছে। তাই কোনো
সিকোয়েন্টের !!a{proof}
\emph{আছে কি না}, যুক্তিবিদ্যার
অন্য পদ্ধতিতেও (যেমন মডেলতত্ত্বে)
তার উত্তর দেওয়া যায়। আর
উভয় ব্যবস্থা ঠিক সেই
সিকোয়েন্টগুলিই !!{prove} করে,
যেগুলি প্রত্যেক !!{structure}-এ
সত্য; ফলে তারা একই সিকোয়েন্ট
প্রমাণ করে।

কিন্তু প্রমাণতত্ত্ব কেবল
কিছু \emph{প্রমাণযোগ্য কি না}
তা নিয়ে নয়, \emph{কীভাবে}
প্রমাণ হয় তাও নিয়ে।
প্রমাণতাত্ত্বিকেরা জিজ্ঞাসা
করেন, “কোনো নির্দিষ্ট বিধি
কি সব সময় বাদ দেওয়া যায়?”,
“এই বিধি যোগ করলে কি নতুন
সিকোয়েন্ট প্রমাণযোগ্য না করেই
ব্যবস্থাটি বাড়ানো যায়?”,
কিংবা “এক ব্যবস্থার !!{proof}s-কে
অন্য ব্যবস্থার !!{proof}s-এ
রূপান্তর করা যায় কি?”
এমন প্রশ্নের উত্তর “হ্যাঁ”
হলে, তা সাধারণত !!{proof}s-এর
জটিলতার ওপর আরোহ করে,
অথবা সমতুল্যভাবে
!!{proof}s-এর গঠনের ওপর আরোহ করে
প্রমাণ করা হয়। এতে শুধু
ফলটির প্রমাণই মেলে না
(যেমন \Log{G1c} কোনো
সিকোয়েন্ট প্রমাণ করলে
\Log{G3c}-ও সেটি প্রমাণ করে),
একটি পদ্ধতিও মেলে
(যেমন \Log{G1c}-র !!{proof}s-কে
\Log{G3c}-র !!{proof}s-এ
রূপান্তর করার পদ্ধতি)।

প্রমাণতত্ত্বের একটি কেন্দ্রীয়
ফল হলো \emph{কর্তন-বর্জন
উপপাদ্য}। এটি দেখায়,
সিকোয়েন্ট ব্যবস্থাগুলিতে
এই কর্তন-বিধি যোগ করলেও
\[
\Axiom$\Gamma \fCenter \Delta, !A$
\Axiom$!A, \Pi \fCenter \Lambda$
\RightLabel{\Cut}
\BinaryInf$\Gamma, \Pi \fCenter \Delta, \Lambda$
\DisplayProof
\]
কোন সিকোয়েন্টগুলি
!!{provable}, তা বদলায় না।
অধিকন্তু প্রমাণটি এমন
একটি পদ্ধতি দেয়, যা
\Log{G1c} (অথবা \Log{G3c})-র
বিধির সঙ্গে \Cut{} ব্যবহার
করা যেকোনো !!{proof}-কে
শুধু \Log{G1c} (অথবা
\Log{G3c})-র বিধি ব্যবহার
করা প্রমাণে রূপান্তর করে।
অর্থাৎ পদ্ধতিটি \Cut{}
ব্যবহারকারী !!{proof}s থেকে
\Cut{} বিধিটিকে “বর্জন” করে;
তাই এই নাম।

\end{document}
